(a) Pullback preserves local trivializations, so $f^*\mathcal L$ is invertible whenever $\mathcal L$ is. The natural isomorphisms $f^*(\mathcal L\otimes\mathcal M)\cong f^*\mathcal L\otimes f^*\mathcal M$ and $f^*\mathcal O_Y\cong\mathcal O_X$ therefore give a homomorphism $f^*:\operatorname{Pic}Y\to\operatorname{Pic}X$.

(b) If $D$ is represented locally by a rational function $g$, then $\mathcal L(D)$ is generated there by $g^{-1}$. Its pullback is generated by $(g\circ f)^{-1}$. At a point $P$ over $Q$, the coefficient of the associated divisor is $v_P(g\circ f)$, which is precisely the coefficient prescribed by the divisor pullback in the text. Thus $f^*\mathcal L(D)\cong\mathcal L(f^*D)$.

(c) For an ambient divisor $D$ having no component containing $X$, its local equations restrict to local equations of the Cartier divisor $D.X$. The same calculation gives $f^*\mathcal L(D)\cong\mathcal L(D.X)$. Every class on $\mathbb P^n$ has such a representative, and (6.16) identifies Cartier and Weil divisor classes on the locally factorial schemes involved. Hence the two homomorphisms agree.
